\subsection{真实输入 \texorpdfstring{$40$}{40} 状态核的 primitive 手术、residue 扭曲谱与 escort 信息几何}\label{subsec:conclusion-realinput40-primitive-surgery-residue-escort-geometry}
本小节把真实输入 \(40\) 状态核中由单原子 Euler 因子 \((1-uz^2)^{-1}\) 所诱导的 primitive 手术、twisted residue 计数、RH 形平方根门槛以及 escort 温度线的信息几何压缩为一组统一的终端命题。其核心结论是：长度 \(2\) 的原子因子可在 \(\zeta\)、ghost 与 primitive 三层被严格反演，但对任何真正的非平凡 twisted 指数严格失明；相应地，所有 RH-sharp 门槛失效都来自 core，而常数层差额则保留为一个严格正的 primitive 质量项。

\begin{theorem}[单原子 primitive 手术在 \texorpdfstring{$\zeta$/ghost/primitive}{zeta/ghost/primitive} 三层上的全函数级可逆性]\label{thm:conclusion-realinput40-atomic-surgery-three-level-invertibility}
\leanverified{paper\_conclusion\_realinput40\_atomic\_surgery\_three\_level\_invertibility}
沿用附录 \ref{app:real-input-40-zeta-u} 的记号。设真实输入 \(40\) 状态核的一参数加权分母满足
\[
\Delta(z,u)=(1-uz^2)\,F(z,u),
\qquad
\zeta(z,u)=\Delta(z,u)^{-1},
\qquad
\zeta_{\mathrm{core}}(z,u):=F(z,u)^{-1},
\]
并记 primitive 生成函数
\[
P(z,u)=\sum_{n\ge 1}p_n(u)z^n,
\qquad
P_{\mathrm{core}}(z,u)=\sum_{n\ge 1}p_n^{\mathrm{core}}(u)z^n,
\]
以及 ghost 展开
\[
\log \zeta(z,u)=\sum_{n\ge 1}\frac{a_n(u)}{n}z^n,
\qquad
\log \zeta_{\mathrm{core}}(z,u)=\sum_{n\ge 1}\frac{a_n^{\mathrm{core}}(u)}{n}z^n.
\]
则单原子剥离在下列三个层面同时严格可逆：
\[
\boxed{\;
P_{\mathrm{core}}(z,u)=P(z,u)-uz^2,\;}
\]
\[
\boxed{\;
\zeta_{\mathrm{core}}(z,u)=(1-uz^2)\,\zeta(z,u),\;}
\]
\[
\boxed{\;
a_n^{\mathrm{core}}(u)=a_n(u)-a_n^{\mathrm{cycle}}(u),\qquad
a_n^{\mathrm{cycle}}(u)=
\begin{cases}
2u^{n/2},&2\mid n,\\
0,&2\nmid n,
\end{cases}}
\]
并且相应 primitive 原子项满足
\[
\boxed{\;
p_n^{\mathrm{cycle}}(u)=
\begin{cases}
u,& n=2,\\
0,& n\neq 2.
\end{cases}}
\]
因此，单原子 surgery 不仅给出向前的剥离，而且定义了一个在 \(\zeta\)、ghost 与 primitive 三层完全兼容的信息无损反演算子。
\end{theorem}
\begin{proof}
由命题 \ref{prop:atomic-2-orbit-split-logM}，\(\zeta\) 层与 primitive 层已经满足
\[
\zeta(z,u)=(1-uz^2)^{-1}\zeta_{\mathrm{core}}(z,u),
\qquad
P(z,u)=uz^2+P_{\mathrm{core}}(z,u).
\]
因此前两条反演公式立即成立。

再由 Witt--ghost 关系，
\[
\log \zeta(z,u)=\sum_{n\ge 1}\frac{a_n(u)}{n}z^n,
\qquad
\log \zeta_{\mathrm{core}}(z,u)=\sum_{n\ge 1}\frac{a_n^{\mathrm{core}}(u)}{n}z^n.
\]
对
\(
\zeta(z,u)=(1-uz^2)^{-1}\zeta_{\mathrm{core}}(z,u)
\)
取对数得
\[
\log \zeta(z,u)
=
-\log(1-uz^2)+\log \zeta_{\mathrm{core}}(z,u)
=
\sum_{m\ge 1}\frac{u^m}{m}z^{2m}
+\log \zeta_{\mathrm{core}}(z,u).
\]
比较 \(z^n\) 系数便得
\[
a_n(u)-a_n^{\mathrm{core}}(u)=
\begin{cases}
2u^{n/2},&2\mid n,\\
0,&2\nmid n.
\end{cases}
\]
这与注 \ref{rem:real-input-40-single-euler} 中的 \(a_n^{\mathrm{cycle}}(u)\) 完全一致。再对 ghost 与 primitive 作 Möbius 反演，即得
\[
p_n^{\mathrm{cycle}}(u)=u\,\mathbf 1_{\{n=2\}}.
\]
于是三个层面的反演彼此一致且唯一。证毕。
\end{proof}

\begin{corollary}[非平凡 twisted 指数的 core 正交性与 RH-sharp 门槛等价]\label{cor:conclusion-realinput40-core-orthogonal-twisted-exponent}
沿用定理 \ref{thm:conclusion-realinput40-atomic-surgery-three-level-invertibility} 的记号。对任意模数 \(m\ge 2\) 与 \(1\le j\le m-1\)，令
\[
u=\omega_m^j,\qquad
\rho_{m,j}^{\mathrm{full}}
:=
\frac{1}{\min\{|z|:\ \Delta(z,u)=0\}},
\qquad
\rho_{m,j}^{\mathrm{core}}
:=
\frac{1}{\min\{|z|:\ F(z,u)=0\}},
\]
并记
\[
\lambda:=\rho(M(1)),\qquad
\theta_{m,j}^{\mathrm{full}}
:=
\frac{\log \rho_{m,j}^{\mathrm{full}}}{\log \lambda},
\qquad
\theta_{m,j}^{\mathrm{core}}
:=
\frac{\log \rho_{m,j}^{\mathrm{core}}}{\log \lambda}.
\]
则有精确关系
\[
\boxed{\;
\rho_{m,j}^{\mathrm{full}}=\max\{1,\rho_{m,j}^{\mathrm{core}}\},\qquad
\theta_{m,j}^{\mathrm{full}}=\max\{0,\theta_{m,j}^{\mathrm{core}}\}.\;}
\]
特别地，
\[
\theta_{m,j}^{\mathrm{full}}>\frac12
\quad\Longleftrightarrow\quad
\theta_{m,j}^{\mathrm{core}}>\frac12.
\]
因此，任何 RH-sharp 门槛失效都严格来自 core 因子 \(F(z,u)\)；被剥离的原子 Euler 因子 \((1-uz^2)^{-1}\) 绝不会制造新的正 twisted 指数。
\end{corollary}
\begin{proof}
前两条盒装恒等式正是定理 \ref{thm:xi-time-part57b-atomic-factor-zero-twisted-exponent} 的结论；而平方根门槛的等价性即推论 \ref{cor:xi-time-part57b-core-full-square-root-barrier-equivalence}。证毕。
\end{proof}

\begin{theorem}[twisted residue 计数的有限 Fourier 全恢复与单轨 delta 支撑]\label{thm:conclusion-realinput40-residue-fourier-recovery-delta-trajectory}
固定模数 \(m\ge 2\)，并记 \(\omega_m:=e^{2\pi i/m}\)。对 \(r\in \ZZ/m\ZZ\) 与 \(n\ge 1\)，定义 periodic 与 primitive 的 energy-residue 计数
\[
A_{m,r}^{\mathrm{full}}(n)
:=
\#\{\gamma:\ |\gamma|=n,\ E(\gamma)\equiv r\pmod m\},
\]
\[
P_{m,r}^{\mathrm{full}}(n)
:=
\#\{\mathfrak p:\ |\mathfrak p|=n,\ E(\mathfrak p)\equiv r\pmod m\}.
\]
相应地，对 core 部分定义 \(A_{m,r}^{\mathrm{core}}(n)\) 与 \(P_{m,r}^{\mathrm{core}}(n)\)。则：
\begin{enumerate}
  \item 完整的 twisted 数据 \(\{a_n(\omega_m^j)\}_{j=0}^{m-1}\) 与 \(\{p_n(\omega_m^j)\}_{j=0}^{m-1}\) 可精确恢复全部 residue 计数：
  \[
  \boxed{\;
  A_{m,r}^{\mathrm{full}}(n)
  =
  \frac1m\sum_{j=0}^{m-1}a_n(\omega_m^j)\,\omega_m^{-jr},\qquad
  P_{m,r}^{\mathrm{full}}(n)
  =
  \frac1m\sum_{j=0}^{m-1}p_n(\omega_m^j)\,\omega_m^{-jr}.\;}
  \]
  \item 原子部分在 residue 图像中只贡献单轨 delta 支撑：
  \[
  \boxed{\;
  A_{m,r}^{\mathrm{full}}(n)
  =
  A_{m,r}^{\mathrm{core}}(n)
  +
  2\,\mathbf 1_{\{2\mid n\}}\mathbf 1_{\{r\equiv n/2\;(\mathrm{mod}\,m)\}},\;}
  \]
  \[
  \boxed{\;
  P_{m,r}^{\mathrm{full}}(n)
  =
  P_{m,r}^{\mathrm{core}}(n)
  +
  \mathbf 1_{\{n=2\}}\mathbf 1_{\{r\equiv 1\;(\mathrm{mod}\,m)\}}.\;}
  \]
\end{enumerate}
因此，periodic 层上的原子贡献只沿 \(r\equiv n/2\pmod m\) 这条直线轨道传播，而 primitive 层上的原子贡献仅支撑在唯一点 \((n,r)=(2,1)\)。
\end{theorem}
\begin{proof}
第一条是有限循环群上的离散 Fourier 反演。对第二条，由定理 \ref{thm:conclusion-realinput40-atomic-surgery-three-level-invertibility}，
\[
a_n(\omega_m^j)-a_n^{\mathrm{core}}(\omega_m^j)
=
2\,\omega_m^{j n/2}\mathbf 1_{\{2\mid n\}},
\]
\[
p_n(\omega_m^j)-p_n^{\mathrm{core}}(\omega_m^j)
=
\omega_m^j\,\mathbf 1_{\{n=2\}}.
\]
把这两式分别代回 Fourier 反演，并利用单位根正交关系
\[
\frac1m\sum_{j=0}^{m-1}\omega_m^{j(\ell-r)}
=
\begin{cases}
1,& \ell\equiv r\pmod m,\\
0,& \ell\not\equiv r\pmod m,
\end{cases}
\]
即可得到两条 delta 支撑公式。等价表述亦已由定理 \ref{thm:xi-atomic-witt-energy-residue-l2-law} 与定理 \ref{thm:xi-root-unity-filter-single-class-injection} 给出。证毕。
\end{proof}

\begin{theorem}[periodic residue 概率律的移动 \texorpdfstring{$\delta$}{delta} 凸分解]\label{thm:conclusion-realinput40-periodic-residue-moving-delta-mixture}
\leanverified{paper\_conclusion\_realinput40\_periodic\_residue\_moving\_delta\_mixture}
固定整数 \(m\ge 2\)。对 \(n\ge 1\) 记
\[
A_n^{\mathrm{full}}
:=
a_n(1)
=
\sum_{r\in\ZZ/m\ZZ}A_{m,r}^{\mathrm{full}}(n),
\qquad
A_n^{\mathrm{core}}
:=
a_n^{\mathrm{core}}(1)
=
\sum_{r\in\ZZ/m\ZZ}A_{m,r}^{\mathrm{core}}(n).
\]
若相应总质量为正，则定义 residue 概率向量
\[
\mu_{m,\bullet}^{\mathrm{full}}(n)
:=
\frac{1}{A_n^{\mathrm{full}}}
\bigl(A_{m,r}^{\mathrm{full}}(n)\bigr)_{r\in\ZZ/m\ZZ},
\qquad
\mu_{m,\bullet}^{\mathrm{core}}(n)
:=
\frac{1}{A_n^{\mathrm{core}}}
\bigl(A_{m,r}^{\mathrm{core}}(n)\bigr)_{r\in\ZZ/m\ZZ}.
\]
则有下列严格分解：
\begin{enumerate}
  \item 若 \(n\) 为奇数，则
  \[
  \mu_{m,\bullet}^{\mathrm{full}}(n)=\mu_{m,\bullet}^{\mathrm{core}}(n).
  \]
  \item 若 \(n=2\ell\) 且 \(A_{2\ell}^{\mathrm{core}}>0\)，则
  \[
  \boxed{\;
  \mu_{m,\bullet}^{\mathrm{full}}(2\ell)
  =
  \frac{A_{2\ell}^{\mathrm{core}}}{A_{2\ell}^{\mathrm{core}}+2}\,
  \mu_{m,\bullet}^{\mathrm{core}}(2\ell)
  +
  \frac{2}{A_{2\ell}^{\mathrm{core}}+2}\,
  \delta_{\ell \bmod m}.\;}
  \]
\end{enumerate}
特别地，当 \(A_{2\ell}^{\mathrm{core}}>0\) 时，full periodic residue law 是 core law 与移动 Dirac 质量 \(\delta_{\ell\bmod m}\) 的严格凸组合；若 \(A_{2\ell}^{\mathrm{core}}=0\)，则
\[
\mu_{m,\bullet}^{\mathrm{full}}(2\ell)=\delta_{\ell\bmod m}.
\]
\end{theorem}
\begin{proof}
奇数情形由定理 \ref{thm:conclusion-realinput40-residue-fourier-recovery-delta-trajectory} 的 periodic residue 修正公式立即得到。若 \(n=2\ell\)，同一定理给出
\[
A_{m,r}^{\mathrm{full}}(2\ell)
=
A_{m,r}^{\mathrm{core}}(2\ell)
+2\,\mathbf 1_{\{r\equiv \ell\;(\mathrm{mod}\,m)\}}.
\]
对 \(r\in\ZZ/m\ZZ\) 求和便得
\[
A_{2\ell}^{\mathrm{full}}=A_{2\ell}^{\mathrm{core}}+2.
\]
若 \(A_{2\ell}^{\mathrm{core}}>0\)，则对上式逐坐标除以 \(A_{2\ell}^{\mathrm{full}}\) 即得到所示分解；若 \(A_{2\ell}^{\mathrm{core}}=0\)，则上式直接退化为
\[
A_{m,r}^{\mathrm{full}}(2\ell)=2\,\mathbf 1_{\{r\equiv \ell\;(\mathrm{mod}\,m)\}},
\]
故 \(\mu_{m,\bullet}^{\mathrm{full}}(2\ell)=\delta_{\ell\bmod m}\)。证毕。
\end{proof}

\begin{theorem}[中心化 residue 扰动的 Fourier 刚性与显式范数]\label{thm:conclusion-realinput40-centered-residue-fourier-rigidity}
固定整数 \(m\ge 2\)。对偶长度 \(n=2\ell\) 定义
\[
\Delta_{m,r}^{(\ell)}
:=
A_{m,r}^{\mathrm{full}}(2\ell)-A_{m,r}^{\mathrm{core}}(2\ell)
=
2\,\mathbf 1_{\{r\equiv \ell\;(\mathrm{mod}\,m)\}},
\]
\[
\widetilde\Delta_{m,r}^{(\ell)}
:=
\Delta_{m,r}^{(\ell)}-\frac{2}{m}.
\]
若 \(f=(f_r)_{r\in\ZZ/m\ZZ}\in\CC^m\)，定义其离散 Fourier 变换为
\[
\widehat f_{m,j}:=\sum_{r=0}^{m-1}f_r\,\omega_m^{jr}
\qquad(0\le j\le m-1),
\]
并记 \(\mathbf 1:=(1,\dots,1)\in\CC^m\)。则：
\begin{enumerate}
  \item \(\widetilde\Delta_{m,\bullet}^{(\ell)}\) 落在零均值超平面
  \[
  \mathcal H_m:=\left\{x\in\CC^m:\ \sum_{r=0}^{m-1}x_r=0\right\},
  \]
  且其 Fourier 变换满足
  \[
  \boxed{\;
  \widehat{\widetilde\Delta}_{m,j}^{(\ell)}
  =
  \begin{cases}
  0,& j=0,\\
  2\,\omega_m^{j\ell},& 1\le j\le m-1.
  \end{cases}}
  \]
  因而全部非平凡 Fourier 模式具有相同模长 \(2\)。
  \item 其 \(\ell^1,\ell^2,\ell^\infty\) 范数具有闭式
  \[
  \boxed{\;
  \bigl\|\widetilde\Delta_{m,\bullet}^{(\ell)}\bigr\|_{\ell^1}
  =
  4\left(1-\frac1m\right),\qquad
  \bigl\|\widetilde\Delta_{m,\bullet}^{(\ell)}\bigr\|_{\ell^2}^2
  =
  4\left(1-\frac1m\right),\qquad
  \bigl\|\widetilde\Delta_{m,\bullet}^{(\ell)}\bigr\|_{\ell^\infty}
  =
  2\left(1-\frac1m\right).\;}
  \]
  \item 对任意满足 \(Q\mathbf 1=0\) 的正交投影 \(Q:\CC^m\to\CC^m\)，都有
  \[
  \boxed{\;
  \bigl\|Q\widetilde\Delta_{m,\bullet}^{(\ell)}\bigr\|_{\ell^2}
  \le
  2\sqrt{1-\frac1m}.\;}
  \]
\end{enumerate}
因此，单原子在中心化 residue 空间中始终由同一参考剖面的循环平移得到：其中一个坐标等于 \(2-\frac{2}{m}\)，其余坐标都等于 \(-\frac{2}{m}\)；在实空间上它表现为单点 spike 减去均匀背景，在 Fourier 空间上则对全部非平凡频率赋予等模振幅。
\end{theorem}
\begin{proof}
由定理 \ref{thm:conclusion-realinput40-residue-fourier-recovery-delta-trajectory}，
\[
\Delta_{m,r}^{(\ell)}
=
2\,\mathbf 1_{\{r\equiv \ell\;(\mathrm{mod}\,m)\}},
\qquad
\widetilde\Delta_{m,r}^{(\ell)}
=
2\,\mathbf 1_{\{r\equiv \ell\;(\mathrm{mod}\,m)\}}-\frac{2}{m}.
\]
因此
\[
\sum_{r=0}^{m-1}\widetilde\Delta_{m,r}^{(\ell)}=2-2=0,
\]
故 \(\widetilde\Delta_{m,\bullet}^{(\ell)}\in \mathcal H_m\)。进一步，
\[
\widehat{\widetilde\Delta}_{m,j}^{(\ell)}
=
2\,\omega_m^{j\ell}
-\frac{2}{m}\sum_{r=0}^{m-1}\omega_m^{jr}.
\]
当 \(j=0\) 时，第二项等于 \(2\)，故结果为 \(0\)；当 \(1\le j\le m-1\) 时，几何级数 \(\sum_{r=0}^{m-1}\omega_m^{jr}=0\)，从而得到所示 Fourier 公式。

再看范数。向量 \(\widetilde\Delta_{m,\bullet}^{(\ell)}\) 恰有一个坐标等于 \(2-\frac{2}{m}\)，其余 \(m-1\) 个坐标都等于 \(-\frac{2}{m}\)。因此
\[
\bigl\|\widetilde\Delta_{m,\bullet}^{(\ell)}\bigr\|_{\ell^1}
=
\left(2-\frac{2}{m}\right)+(m-1)\frac{2}{m}
=
4\left(1-\frac1m\right),
\]
\[
\bigl\|\widetilde\Delta_{m,\bullet}^{(\ell)}\bigr\|_{\ell^2}^2
=
\left(2-\frac{2}{m}\right)^2
+(m-1)\left(\frac{2}{m}\right)^2
=
4\left(1-\frac1m\right),
\]
\[
\bigl\|\widetilde\Delta_{m,\bullet}^{(\ell)}\bigr\|_{\ell^\infty}
=
2-\frac{2}{m}.
\]
最后，正交投影在 \(\ell^2\) 范数下是压缩映射，故
\[
\bigl\|Q\widetilde\Delta_{m,\bullet}^{(\ell)}\bigr\|_{\ell^2}
\le
\bigl\|\widetilde\Delta_{m,\bullet}^{(\ell)}\bigr\|_{\ell^2}
=
2\sqrt{1-\frac1m}.
\]
证毕。
\end{proof}

\begin{theorem}[primitive residue 生成函数的 monomial 刚性]\label{thm:conclusion-realinput40-primitive-residue-monomial-rigidity}
固定整数 \(m\ge 2\) 与剩余类 \(r\in\ZZ/m\ZZ\)。定义 primitive residue 生成函数
\[
\mathcal P_{m,r}^{\mathrm{full}}(z)
:=
\sum_{n\ge 1}P_{m,r}^{\mathrm{full}}(n)\,z^n,
\qquad
\mathcal P_{m,r}^{\mathrm{core}}(z)
:=
\sum_{n\ge 1}P_{m,r}^{\mathrm{core}}(n)\,z^n.
\]
则有严格恒等式
\[
\boxed{\;
\mathcal P_{m,r}^{\mathrm{full}}(z)
=
\mathcal P_{m,r}^{\mathrm{core}}(z)
+z^2\,\mathbf 1_{\{r\equiv 1\;(\mathrm{mod}\,m)\}}.\;}
\]
因此，\(\mathcal P_{m,r}^{\mathrm{full}}\) 与 \(\mathcal P_{m,r}^{\mathrm{core}}\) 具有相同的收敛半径、相同的全部非多项式奇点，以及相同的系数指数增长率
\[
\limsup_{n\to\infty}\bigl(P_{m,r}^{\mathrm{full}}(n)\bigr)^{1/n}
=
\limsup_{n\to\infty}\bigl(P_{m,r}^{\mathrm{core}}(n)\bigr)^{1/n}.
\]
特别地，凡由 residue-resolved primitive 生成函数的主导奇点、收敛半径或系数指数所定义的 primitive 扭曲指数与 RH-sharp 阈值，在 full 与 core 之间完全一致。
\end{theorem}
\begin{proof}
由定理 \ref{thm:conclusion-realinput40-residue-fourier-recovery-delta-trajectory}，
\[
P_{m,r}^{\mathrm{full}}(n)-P_{m,r}^{\mathrm{core}}(n)
=
\mathbf 1_{\{n=2\}}\mathbf 1_{\{r\equiv 1\;(\mathrm{mod}\,m)\}}.
\]
对 \(n\ge 1\) 求和即可得到
\[
\mathcal P_{m,r}^{\mathrm{full}}(z)
=
\mathcal P_{m,r}^{\mathrm{core}}(z)
+z^2\,\mathbf 1_{\{r\equiv 1\;(\mathrm{mod}\,m)\}}.
\]
两者之差只是一个次数 \(2\) 的多项式，因此其收敛半径与全部非多项式奇点完全一致。再由 Cauchy--Hadamard 公式，系数指数增长率只由收敛半径决定，故上式中的两侧极限上确界相同。最后一句只是这一事实对由主导奇点或指数增长率定义之阈值量的直接重述。证毕。
\end{proof}

\begin{corollary}[periodic residue 概率律的精确 \texorpdfstring{$W_1$}{W1} 位移公式]\label{cor:conclusion-realinput40-periodic-residue-w1-exact}
固定整数 \(m\ge 2\) 与 \(\ell\ge 1\)，并假设 \(A_{2\ell}^{\mathrm{core}}>0\)。在 \(\ZZ/m\ZZ\) 上赋予循环距离
\[
d_m(r,s):=\min_{k\in\ZZ}|r-s+km|.
\]
记
\[
\varepsilon_{m,\ell}:=\frac{2}{A_{2\ell}^{\mathrm{core}}+2}.
\]
则由定理 \ref{thm:conclusion-realinput40-periodic-residue-moving-delta-mixture}，
\[
\mu_{m,\bullet}^{\mathrm{full}}(2\ell)
=
\bigl(1-\varepsilon_{m,\ell}\bigr)\mu_{m,\bullet}^{\mathrm{core}}(2\ell)
+\varepsilon_{m,\ell}\,\delta_{\ell\bmod m},
\]
并且一阶 Wasserstein 距离满足精确公式
\[
\boxed{\;
W_1\!\left(
\mu_{m,\bullet}^{\mathrm{full}}(2\ell),
\mu_{m,\bullet}^{\mathrm{core}}(2\ell)
\right)
=
\varepsilon_{m,\ell}
\sum_{r\in\ZZ/m\ZZ}
d_m(r,\ell)\,\mu_{m,r}^{\mathrm{core}}(2\ell).\;}
\]
\end{corollary}
\begin{proof}
记
\[
\mu:=\mu_{m,\bullet}^{\mathrm{core}}(2\ell),
\qquad
\nu:=\mu_{m,\bullet}^{\mathrm{full}}(2\ell),
\qquad
\varepsilon:=\varepsilon_{m,\ell}.
\]
则
\[
\nu=(1-\varepsilon)\mu+\varepsilon\delta_{\ell\bmod m}.
\]
因此对每个 \(r\neq \ell\bmod m\)，源分布在 \(r\) 处比目标分布多出恰好 \(\varepsilon\mu_r\) 的质量，而唯一需要补充质量的点是 \(\ell\bmod m\)。故任意运输计划都必须把至少 \(\varepsilon\mu_r\) 的质量从 \(r\) 送往 \(\ell\bmod m\)，从而产生下界
\[
W_1(\nu,\mu)\ge \varepsilon\sum_{r\in\ZZ/m\ZZ}d_m(r,\ell)\,\mu_r.
\]
另一方面，把每个 \(r\) 处的 \(\varepsilon\mu_r\) 质量直接搬运到 \(\ell\bmod m\)，并把剩余的 \((1-\varepsilon)\mu_r\) 留在原位，便得到一条可行耦合，其总代价正好等于右端。于是下界可达，所示公式成立。证毕。
\end{proof}

\begin{theorem}[escort 温度线的 Fisher--Rao 平直化与测地线重参数化]\label{thm:conclusion-escort-fisher-rao-geodesic-reparametrization}
沿用结论 \ref{con:xi-time-part9o-escort-fisher-rao-closed} 的记号。对 \(q\ge 0\)，令
\[
\theta(q):=\frac{1}{1+\varphi^{q+1}},
\qquad
s(q):=2\arcsin\sqrt{\theta(q)}.
\]
则 escort 温度线的极限 Fisher 信息满足
\[
\boxed{\;
\mathcal I(q)=(\log\varphi)^2\,\theta(q)\bigl(1-\theta(q)\bigr).\;}
\]
相应极限 Fisher--Rao 度量可写为
\[
\boxed{\;
\mathrm ds^2=\mathcal I(q)\,\mathrm dq^2=\bigl(\mathrm d s(q)\bigr)^2.\;}
\]
并且 \(q\mapsto s(q)\) 严格单调，因此 escort 温度线在参数 \(s\) 下成为一条单位速度测地线；其 Fisher--Rao 距离闭式为
\[
\boxed{\;
d_{\mathrm{FR}}(q_0,q_1)
=
\bigl|s(q_1)-s(q_0)\bigr|
=
2\left|
\arcsin\sqrt{\theta(q_1)}-\arcsin\sqrt{\theta(q_0)}
\right|.\;}
\]
\end{theorem}
\begin{proof}
由结论 \ref{con:xi-time-part9o-escort-fisher-rao-closed}，
\[
\mathcal I(q)=(\log\varphi)^2\,\theta(q)\bigl(1-\theta(q)\bigr),
\qquad
\theta'(q)=-(\log\varphi)\,\theta(q)\bigl(1-\theta(q)\bigr).
\]
另一方面，
\[
s'(q)
=
\frac{\theta'(q)}{\sqrt{\theta(q)(1-\theta(q))}}
=
-(\log\varphi)\sqrt{\theta(q)(1-\theta(q))},
\]
故
\[
\bigl(s'(q)\bigr)^2
=
(\log\varphi)^2\,\theta(q)\bigl(1-\theta(q)\bigr)
=
\mathcal I(q).
\]
于是
\[
\mathcal I(q)\,\mathrm dq^2
=
\bigl(\mathrm d s(q)\bigr)^2,
\]
从而 \(s\) 是 Fisher--Rao 度量的平直坐标。又因为 \(\theta'(q)<0\) 对一切 \(q\ge 0\) 成立，故 \(q\mapsto s(q)\) 严格单调。距离公式则由一维平直度量的积分表示立即得到。证毕。
\end{proof}

\begin{theorem}[Abel 常数差额的严格正 primitive 质量解释]\label{thm:conclusion-realinput40-abel-gap-positive-primitive-mass}
沿用附录 \ref{app:real-input-40-zeta-u} 的记号。记
\[
\lambda(u):=\rho(M(u)),\qquad
z_\star(u):=\lambda(u)^{-1},
\]
并定义原子差额函数
\[
\mathcal D_{\mathrm{at}}(u)
:=
\log\mathfrak{M}(u)-\log\mathfrak{M}_{\mathrm{core}}(u).
\]
则对每个 \(u>0\) 都有严格闭式
\[
\boxed{\;
\mathcal D_{\mathrm{at}}(u)=u\,z_\star(u)^2=u\,\lambda(u)^{-2}>0.\;}
\]
特别地，在未扭曲点 \(u=1\) 上，
\[
\mathcal D_{\mathrm{at}}(1)=z_\star^2=\varphi^{-4}=\frac{7-3\sqrt5}{2}.
\]
因此，\(\log\mathfrak{M}(u)-\log\mathfrak{M}_{\mathrm{core}}(u)\) 并非由重整化方案任意产生的差额，而是被剥离的长度 \(2\) primitive 原子在 Perron 归一化下留下的严格正质量。
\end{theorem}
\begin{proof}
对命题 \ref{prop:primitive-orbit-surgery} 取
\[
m=1,\qquad \ell=2,\qquad E=1,
\]
便得到
\[
\log\mathfrak{M}(u)=\log\mathfrak{M}_{\mathrm{core}}(u)+u\,z_\star(u)^2.
\]
这已给出所需闭式。又因 \(u>0\) 时矩阵族 \(M(u)\) 非负，故 Perron 根 \(\lambda(u)>0\)，从而
\[
z_\star(u)=\lambda(u)^{-1}>0.
\]
于是 \(u\,z_\star(u)^2>0\)，严格正性成立。未扭曲点 \(u=1\) 的数值化简则由 \(z_\star=\varphi^{-2}\) 直接得到。证毕。
\end{proof}

\paragraph{40 状态真实输入核的单源临界带、Lucas 素长估值与有限部双机制}
在定理 \ref{thm:conclusion-realinput40-atomic-surgery-three-level-invertibility} 已把单原子 surgery 的 \(\zeta\)/ghost/primitive 三层反演固定、而定理 \ref{thm:conclusion-realinput40-abel-gap-positive-primitive-mass} 已把原子差额的 Abel 质量闭式锁定之后，还可把真实输入 \(40\) 状态核中由 primitive 轨道谱拆分、Dirichlet 完成截面、Lucas 型素长闭式与 Perron 有限部接口共同诱导的解析几何进一步压缩为下列五条后果。以下固定
\[
\varphi=\frac{1+\sqrt5}{2},
\qquad
F:=\begin{pmatrix}1&1\\[2pt]1&0\end{pmatrix},
\qquad
\lambda:=\rho(M)=\varphi^2,
\qquad
z_\star:=\lambda^{-1}=\varphi^{-2},
\]
并记
\[
\widehat{\zeta}_M(s):=\zeta_M(\lambda^{-s}).
\]

\begin{theorem}[开临界带的单源分解]\label{thm:conclusion-realinput40-open-strip-single-source}
\leanverified{paper\_conclusion\_realinput40\_open\_strip\_single\_source}
记
\[
H(s):=
\frac{\zeta_{F\otimes F}(\lambda^{-s})}
{(1-\lambda^{-s})^2(1+\lambda^{-s})(1-\varphi^{-1}\lambda^{-s})}.
\]
则 \(H\) 在开条带 \(0<\Re(s)<1\) 内全纯且无零，并且
\[
\widehat{\zeta}_M(s)=\frac{H(s)}{1+\varphi^{1-2s}}.
\]
特别地，\(\widehat{\zeta}_M\) 在开临界带中的全部极点都由单一因子 \(1+\varphi^{1-2s}\) 产生。
\end{theorem}
\begin{proof}
由命题 \ref{prop:real-input-40-fib-tensor}，
\[
\zeta_M(z)=\frac{1}{(1-z)^2(1+z)}\,
\zeta_{F\otimes F}(z)\,\zeta_{-F}(z),
\qquad
\zeta_{-F}(z)=\frac{1}{1+z-z^2}.
\]
又
\[
1+z-z^2=(1-\varphi^{-1}z)(1+\varphi z),
\]
故代入 \(z=\lambda^{-s}\) 得
\[
\widehat{\zeta}_M(s)
=
\frac{\zeta_{F\otimes F}(\lambda^{-s})}
{(1-\lambda^{-s})^2(1+\lambda^{-s})(1-\varphi^{-1}\lambda^{-s})(1+\varphi\lambda^{-s})}
=
\frac{H(s)}{1+\varphi^{1-2s}}.
\]
其中 \(1+\varphi\lambda^{-s}=1+\varphi^{1-2s}\)。另一方面，
\[
\zeta_{F\otimes F}(z)=\frac{1}{(1-\lambda z)(1+z)^2(1-\lambda^{-1}z)},
\]
故 \(H\) 的可能奇点只能落在 \(\Re(s)\in\{1,0,-\tfrac12,-1\}\) 上，因而在 \(0<\Re(s)<1\) 内全纯。又 \(\zeta_{F\otimes F}\) 与分母各因子都不在该开条带内取零，故 \(H\) 在该条带内无零。证毕。
\end{proof}

\begin{theorem}[临界带极点晶格的等距同留数性]\label{thm:conclusion-realinput40-open-strip-pole-lattice-equal-residue}
设
\[
s_k:=\frac12+\frac{(2k+1)\pi i}{\log\lambda},
\qquad
k\in\ZZ.
\]
则 \(\widehat{\zeta}_M\) 在开临界带中的全部极点恰为 \(\{s_k\}_{k\in\ZZ}\)；这些极点全为单极点，并满足
\[
\operatorname*{Res}_{s=s_k}\widehat{\zeta}_M(s)
=
\frac{H(s_0)}{\log\lambda}
\qquad
(k\in\ZZ).
\]
\end{theorem}
\begin{proof}
由定理 \ref{thm:conclusion-realinput40-open-strip-single-source}，开临界带内的极点恰由
\[
1+\varphi^{1-2s}=0
\]
决定。该方程等价于
\[
(1-2s)\log\varphi=(2k+1)\pi i,
\]
亦即
\[
s=s_k=\frac12+\frac{(2k+1)\pi i}{\log\lambda}.
\]
再对分母求导：
\[
\frac{\mathrm d}{\mathrm ds}\bigl(1+\varphi^{1-2s}\bigr)
=-\log\lambda\cdot \varphi^{1-2s}.
\]
在 \(s=s_k\) 处有 \(\varphi^{1-2s_k}=-1\)，故该导数等于 \(\log\lambda\neq 0\)，从而所有极点均为单极点。最后，\(H\) 只依赖于 \(\lambda^{-s}\)，因而具有虚周期 \(2\pi i/\log\lambda\)；而
\[
s_{k+1}-s_k=\frac{2\pi i}{\log\lambda},
\]
所以 \(H(s_k)=H(s_0)\) 对一切 \(k\) 成立。留数公式遂化为
\[
\operatorname*{Res}_{s=s_k}\widehat{\zeta}_M(s)
=
\frac{H(s_k)}{\log\lambda}
=
\frac{H(s_0)}{\log\lambda}.
\]
证毕。
\end{proof}

\begin{theorem}[primitive 平方根异常的奇支撑]\label{thm:conclusion-realinput40-primitive-sqrt-anomaly-odd-support}
记 \(p_n:=p_n(M)\)。则有统一渐近
\[
p_n
=
\frac{\lambda^n}{n}
-\mathbf 1_{\{n\ \mathrm{odd}\}}\frac{\lambda^{n/2}}{n}
+O\!\left(\frac{\lambda^{n/3}}{n}\right).
\]
等价地，
\[
p_n
=
\frac{\lambda^n}{n}
+\frac{\bigl((-1)^n-\mathbf 1_{\{2\mid n\}}\bigr)\lambda^{n/2}}{n}
+O\!\left(\frac{\lambda^{n/3}}{n}\right).
\]
特别地，\(\lambda^{n/2}\) 级别的平方根异常完全支撑在奇长度上；在偶长度上，M\"obius 反演中的 \(d=1\) 与 \(d=2\) 两项于该尺度发生精确抵消。
\end{theorem}
\begin{proof}
命题 \ref{prop:finite-rh-40} 已给出：真实输入 \(40\) 状态核在 \(|\mu|=\sqrt{\lambda}\) 上恰有唯一非 Perron 模态 \(-\sqrt{\lambda}\)，而其余特征值严格落在 \(|\mu|<\sqrt{\lambda}\) 内。因此定理 \ref{thm:primitive-odd-even-sqrt-cancellation} 可直接应用，并在 \(\rho=\lambda\) 时给出
\[
p_n
=
\frac{\lambda^n}{n}
-\mathbf 1_{\{n\ \mathrm{odd}\}}\frac{\lambda^{n/2}}{n}
+O\!\left(\frac{\lambda^{n/3}}{n}\right).
\]
第二个写法只是把奇偶指示函数作代数重写。最后一句正是同一定理中关于 \(d=1\) 与 \(d=2\) 精确消项机制的直接重述。证毕。
\end{proof}

\begin{theorem}[素长 primitive 层的 Lucas--Wieferich 估值律]\label{thm:conclusion-realinput40-prime-length-lucas-wieferich-valuation}
设 \(p\) 为奇素数，记
\[
L_p:=\Tr(F^p),
\qquad
p_p:=p_p(M),
\qquad
\nu_p:=v_p(L_p-1).
\]
则有精确公式
\[
p_p=\frac{L_p(L_p-1)}{p},
\qquad
L_p\equiv 1 \pmod p,
\qquad
v_p(p_p)=\nu_p-1.
\]
因而对任意 \(k\ge 0\)，
\[
p^k\mid p_p
\iff
p^{k+1}\mid (L_p-1),
\qquad
p_p\equiv \frac{L_p-1}{p}\pmod p.
\]
\end{theorem}
\begin{proof}
推论 \ref{cor:real-input-40-prime-length-product} 已给出
\[
p_p=\frac{L_p(L_p-1)}{p}.
\]
另一方面，引理 \ref{lem:pom-lucas-modp} 表明 \(L_p\equiv 1\pmod p\)，故 \(v_p(L_p)=0\)。于是
\[
v_p(p_p)=v_p(L_p)+v_p(L_p-1)-1=\nu_p-1.
\]
这立即推出
\[
p^k\mid p_p
\iff
v_p(p_p)\ge k
\iff
\nu_p\ge k+1
\iff
p^{k+1}\mid(L_p-1).
\]
又因 \(L_p\equiv 1\pmod p\)，故在模 \(p\) 意义下 \(L_p\) 为单位并且同余于 \(1\)，从而
\[
p_p=\frac{L_p(L_p-1)}{p}\equiv \frac{L_p-1}{p}\pmod p.
\]
证毕。
\end{proof}

\begin{theorem}[Perron 极留数的局域手术律与 Abel 有限部的三因子加法律]\label{thm:conclusion-realinput40-perron-residue-local-surgery-finitepart-additivity}
将 \(\zeta_M\) 写为
\[
\zeta_M(z)=\zeta_{F\otimes F}(z)\,R_{\mathrm{reg}}(z),
\qquad
R_{\mathrm{reg}}(z):=\frac{\zeta_{-F}(z)}{(1-z)^2(1+z)}
=\frac{1}{(1+z-z^2)(1-z)^2(1+z)}.
\]
则有：
\begin{enumerate}
  \item \(R_{\mathrm{reg}}\) 在 \(z=z_\star\) 处解析且非零，并满足显式闭式
  \[
  R_{\mathrm{reg}}(z_\star)=\frac{15+7\sqrt5}{20}.
  \]
  \item 若记
  \[
  C_M:=\lim_{z\to z_\star}(1-\lambda z)\zeta_M(z),
  \qquad
  C_{F\otimes F}:=\lim_{z\to z_\star}(1-\lambda z)\zeta_{F\otimes F}(z),
  \]
  则
  \[
  C_M=R_{\mathrm{reg}}(z_\star)\,C_{F\otimes F}.
  \]
  \item 若按注 \ref{rem:real-input-40-logM-artin-factorization} 记 \(\log\mathfrak{M}_{\mathrm{even}}\) 为 \(\zeta_{\mathrm{even}}(z)=(1-3z+z^2)^{-1}\) 的有限部常数，并定义
  \[
  \log\mathfrak{M}_{\mathrm{twist}}:=K_{-F},
  \qquad
  \log\mathfrak{M}_{\mathrm{triv}}:=P_{\mathrm{triv}}(z_\star),
  \]
  则 Abel 有限部满足
  \[
  \log\mathfrak{M}
  =
  \log\mathfrak{M}_{\mathrm{even}}
  +\log\mathfrak{M}_{\mathrm{twist}}
  +\log\mathfrak{M}_{\mathrm{triv}}.
  \]
\end{enumerate}
因此，Perron 主极点的留数修正满足纯局域的乘法手术律，而 Abel 有限部则保持 even/twist/triv 三因子的全局加法律。
\end{theorem}
\begin{proof}
命题 \ref{prop:real-input-40-fib-tensor} 已给出
\[
\zeta_M(z)=\frac{1}{(1-z)^2(1+z)}\,\zeta_{F\otimes F}(z)\,\zeta_{-F}(z),
\]
故所示 \(R_{\mathrm{reg}}\) 分解成立。由于 \(z_\star=\varphi^{-2}\in(0,1)\) 且 \(\zeta_{-F}\) 的极点位于 \(z=-\varphi^{-1}\) 与 \(z=\varphi\)，故 \(R_{\mathrm{reg}}\) 在 \(z_\star\) 处解析且非零。直接代入并化简得
\[
R_{\mathrm{reg}}(z_\star)
=
\frac{1}{(1+\varphi^{-2}-\varphi^{-4})(1-\varphi^{-2})^2(1+\varphi^{-2})}
=
\frac{15+7\sqrt5}{20}.
\]
于是第二条只需把 \(\zeta_M=\zeta_{F\otimes F}R_{\mathrm{reg}}\) 乘以 \(1-\lambda z\) 后令 \(z\to z_\star\) 即得：
\[
C_M
=
\lim_{z\to z_\star}(1-\lambda z)\zeta_{F\otimes F}(z)\,R_{\mathrm{reg}}(z)
=
R_{\mathrm{reg}}(z_\star)\,C_{F\otimes F}.
\]
最后，注 \ref{rem:real-input-40-logM-artin-factorization} 已明确给出
\[
\log\mathfrak{M}
=
\log\mathfrak{M}_{\mathrm{even}}+K_{-F}+P_{\mathrm{triv}}(z_\star).
\]
将 \(K_{-F}\) 与 \(P_{\mathrm{triv}}(z_\star)\) 重新记作 \(\log\mathfrak{M}_{\mathrm{twist}}\) 与 \(\log\mathfrak{M}_{\mathrm{triv}}\)，便得到第三条。证毕。
\end{proof}

在前述 primitive surgery、ghost 剥离与 Abel 常数差额公式之外，真实输入 \(40\) 状态核还允许把唯一长度 \(2\) 原子与 unitary-slice / Toeplitz 缺陷口径进一步正交化为下列结果。

\begin{theorem}[真实输入 \texorpdfstring{$40$}{40} 状态核中唯一 primitive 二步原子的刚性]\label{thm:conclusion-realinput40-unique-primitive-two-step-atom}
\leanverified{paper\_conclusion\_realinput40\_unique\_primitive\_two\_step\_atom}
沿用附录 \ref{app:real-input-40-zeta-u} 的记号。设
\[
\Delta(z,u)=(1-uz^2)F(z,u),
\qquad
\zeta(z,u)=\Delta(z,u)^{-1},
\]
并记 primitive 生成函数为
\[
P(z,u)=\sum_{n\ge 1}p_n(u)z^n.
\]
则存在唯一分裂
\[
P(z,u)=uz^2+P_{\mathrm{core}}(z,u)
\]
使得
\[
\exp\!\Bigl(\sum_{m\ge 1}\frac{(uz^2)^m}{m}\Bigr)=(1-uz^2)^{-1},
\qquad
\exp\!\Bigl(\sum_{m\ge 1}\frac{P_{\mathrm{core}}(z^m,u^m)}{m}\Bigr)=F(z,u)^{-1}.
\]
特别地：
\begin{enumerate}
  \item \(uz^2\) 是该 family 中唯一被精确剥离出的 primitive 原子；
  \item 任一与因子分解 \(\Delta(z,u)=(1-uz^2)F(z,u)\) 相容的 primitive 分裂，其原子部分必等于 \(uz^2\)；
  \item 该原子恰对应 branch-level 中全部 \((|\tau|,E)=(2,1)\) 的二步 primitive 回路总和。
\end{enumerate}
\end{theorem}
\begin{proof}
由命题 \ref{prop:atomic-2-orbit-split-logM} 与定理 \ref{thm:xi-atomic-witt-cycle-primitive-exact-peeling}，
\[
\zeta(z,u)=(1-uz^2)^{-1}\zeta_{\mathrm{core}}(z,u),
\qquad
P(z,u)=uz^2+P_{\mathrm{core}}(z,u).
\]
这已经给出所述分裂的存在性。对数化后
\[
\log(1-uz^2)^{-1}
=
\sum_{m\ge 1}\frac{u^m z^{2m}}{m}
=
\sum_{m\ge 1}\frac{(uz^2)^m}{m},
\]
故该分裂与 \(\Delta=(1-uz^2)F\) 完全一致。

若另有一组 primitive 分裂
\[
P(z,u)=P_{\mathrm{at}}(z,u)+P_{\mathrm{rem}}(z,u)
\]
满足
\[
\exp\!\Bigl(\sum_{m\ge 1}\frac{P_{\mathrm{at}}(z^m,u^m)}{m}\Bigr)=(1-uz^2)^{-1},
\]
则由 Big-Witt / M\"obius 反演的唯一性，必有
\[
P_{\mathrm{at}}(z,u)=uz^2.
\]
故原子部分唯一。至于 branch-level 解释，命题 \ref{prop:atomic-2-orbit-split-logM} 中对 \((|\tau|,E)=(2,1)\) 的逐项审计已经说明：被剥离的正是全部二步 primitive 回路的总贡献。证毕。
\end{proof}

\begin{theorem}[\texorpdfstring{$2$}{2}-典型 ghost 塔中的二步原子刚性]\label{thm:conclusion-realinput40-two-typical-ghost-irreducibility}
\leanverified{paper\_conclusion\_realinput40\_two\_typical\_ghost\_irreducibility}
设 \(L_{(2,u)}\in \mathrm{TC}_1\) 为定理 \ref{thm:atomic-witt-into-tc1} 由原子数据 \((u,2,1)\) 编译得到的任一迹类模型。则在 \(2\)-典型分量上有：
\begin{enumerate}
  \item 其 ghost 塔由
  \[
  \Bigl\{\Tr\!\bigl(L_{(2,u)}^{2^n}\bigr)\Bigr\}_{n\ge 0}
  \]
  完全给出；
  \item Frobenius--Verschiebung 关系严格实现为
  \[
  \mathsf F_2\!\bigl(\mathsf V_2(L_{(2,u)})\bigr)\simeq L_{(2,u)}^{\oplus 2};
  \]
  \item 在本文允许的 admissible 原子数据类中，不存在只由长度严格小于 \(2\) 的 primitive atoms 生成的有限直和，其全部 \(2\)-幂 ghost 坐标与 \(L_{(2,u)}\) 一致。
\end{enumerate}
因此，长度 \(2\) 原子不仅在 primitive 分裂里唯一，而且在 \(2\)-typical ghost tower 中也保持不可约。
\end{theorem}
\begin{proof}
由定理 \ref{thm:atomic-witt-into-tc1}，原子数据 \((u,2,1)\) 可被编译为迹类算子模型 \(L_{(2,u)}\)，其 ghost 坐标由迹序列 \(\Tr(L_{(2,u)}^n)\) 实现；限制到 \(2\)-幂索引便得到第一条。

第二条是命题 \ref{prop:cyclic-p-typical-frobenius-verschiebung} 在 \(p=2\) 时的直接结论：
\[
\Tr\!\bigl(\mathsf F_2(L)^n\bigr)=\Tr\!\bigl(L^{2n}\bigr),
\qquad
\mathsf F_2\!\bigl(\mathsf V_2(L)\bigr)\simeq L^{\oplus 2}.
\]
取 \(L=L_{(2,u)}\) 即得。

若第三条不成立，则存在只由长度 \(<2\) 的 primitive atoms 构成的 admissible 有限直和 \(L'\)，使
\[
\Tr\!\bigl((L')^{2^n}\bigr)=\Tr\!\bigl(L_{(2,u)}^{2^n}\bigr)
\qquad(\forall n\ge 0).
\]
由 \(2\)-typical Witt 变换的函子性，\(L'\) 与 \(L_{(2,u)}\) 的 \(2\)-典型分量必须一致；但定理 \ref{thm:conclusion-realinput40-unique-primitive-two-step-atom} 与定理 \ref{thm:conclusion-realinput40-root-unity-ghost-complete-primitive-degenerate} 已说明，在本文 admissible primitive 扇区中，唯一非退化原子脉冲正是长度 \(2\) 的 \(uz^2\)，而长度 \(<2\) 的 primitive 支撑不可能产生同一原子类。矛盾。证毕。
\end{proof}

\begin{theorem}[primitive 有限部位移与 Toeplitz 缺陷指数的正交性]\label{thm:conclusion-realinput40-finitepart-diskdefect-orthogonality}
在未扭曲点 \(u=1\) 上，真实输入 \(40\) 状态核的长度 \(2\) primitive 原子同时满足：
\begin{enumerate}
  \item 它引起 Abel 有限部常数的严格非零位移；
  \item 它不产生任何盘内 Blaschke 缺陷，因此不改变 Pontryagin 负平方指数 \(\kappa\)。
\end{enumerate}
更准确地说，若
\[
\Delta(z,u)=(1-uz^2)F(z,u),
\]
则在 \(u=1\) 时
\[
P(z)=z^2+P_{\mathrm{core}}(z),
\qquad
\mathcal D_{\mathrm{at}}(1)=\varphi^{-4}>0,
\]
而因子
\[
1-z^2
\]
的全部零点均位于单位圆上，故对应 Jensen 缺陷为零，从而对应的 Toeplitz / Pontryagin 缺陷指数贡献亦为零。
\end{theorem}
\begin{proof}
由定理 \ref{thm:conclusion-realinput40-unique-primitive-two-step-atom}，
\[
P(z)=z^2+P_{\mathrm{core}}(z)
\]
在 \(u=1\) 时仍保留非平凡长度 \(2\) 原子。再由定理 \ref{thm:conclusion-realinput40-abel-gap-positive-primitive-mass}，
\[
\mathcal D_{\mathrm{at}}(1)=z_\star^2=\varphi^{-4}>0,
\]
故 finite-part / Abel 常数层的位移严格非零。

另一方面，
\[
1-z^2=(1-z)(1+z)
\]
的零点恰为 \(\pm 1\)，全部位于单位圆上。由推论 \ref{cor:app-rh-iff-jensen-defect-zero} 与定理 \ref{thm:xi-selfinversive-jensen-criterion}，这样的 self-inversive 因子对应 Jensen 缺陷恒为零；再由定理 \ref{thm:xi-toeplitz-negative-squares-equals-blaschke-defect}，盘内 Blaschke 缺陷维数与 Pontryagin 负平方指数完全一致，故该因子不贡献任何新的负平方。于是 finite-part displacement 与 disk-defect index 严格正交。证毕。
\end{proof}

\begin{conjecture}[unitary-slice 伪模类的二不变量完备性]\label{conj:conclusion-unitary-slice-pseudomoduli-two-invariants}
在有限缺陷、有限 primitive surgery 的 admissible 完成化 family 类中，离线审计等价类应完全由以下两个不变量决定：
\begin{enumerate}
  \item Pontryagin--Toeplitz 缺陷指数
  \[
  \kappa;
  \]
  \item primitive 有限部位移所对应的 Witt 原子多重集
  \[
  \mathfrak A_{\mathrm{prim}}.
  \]
\end{enumerate}
更具体地说，若两族完成化对象 \(\mathcal F_1,\mathcal F_2\) 满足
\[
\kappa(\mathcal F_1)=\kappa(\mathcal F_2),
\qquad
\mathfrak A_{\mathrm{prim}}(\mathcal F_1)=\mathfrak A_{\mathrm{prim}}(\mathcal F_2),
\]
则它们在 unitary-slice / Toeplitz--PSD / finite-part 三重审计意义下应属于同一伪模类。

该猜想的理由在于：定理 \ref{thm:conclusion-unitary-slice-zero-defect-critical-rigidity} 与定理 \ref{thm:conclusion-toeplitz-failure-complete-geometric-localization} 已说明 \(\kappa\) 完全控制 disk-defect、Toeplitz 负平方与离线零点纤维的几何局域化；而定理 \ref{thm:conclusion-realinput40-unique-primitive-two-step-atom}、定理 \ref{thm:conclusion-realinput40-two-typical-ghost-irreducibility} 与定理 \ref{thm:conclusion-realinput40-finitepart-diskdefect-orthogonality} 又说明 \(\mathfrak A_{\mathrm{prim}}\) 独立控制 primitive surgery、finite-part displacement 与 \(p\)-typical ghost profile，并且与 \(\kappa\) 严格正交。若存在更小的完备不变量，则必须进一步压缩这两条彼此正交的审计轴；现有框架并未显示这种压缩机制。
\end{conjecture}


上述结果把该接口进一步压缩到四层：其一，periodic residue law 在偶长度上只是 core law 与移动 Dirac 质量的凸分解，中心化后的原子扰动在实空间与 Fourier 空间均呈现完全刚性的单轨剖面，而 primitive residue 生成函数则仅多出一项 \(z^2\mathbf 1_{\{r\equiv 1\;(\mathrm{mod}\,m)\}}\)，因此全部 residue-resolved 主导奇点与指数阈值仍由 core 决定；其二，Dirichlet 完成截面在开临界带内退化为由单一负根因子生成的等距同留数极点晶格；其三，primitive 层的平方根异常只支撑在奇长度，而素长算术又被 Lucas--Wieferich 深度完全分级；其四，Perron 主极点的留数修正满足局域乘法手术律，而 Abel 有限部则保持 even/twist/triv 三因子的全局加法律。于是，真实输入 \(40\) 状态核中的 arithmetic surgery、twisted spectrum、Dirichlet 临界几何与 finite-part bookkeeping 便被压缩为同一条可审计的分层结构：指数门槛与非多项式奇性由 core 与单一符号通道共同锁定，而局部概率偏移、极点主部与常数层差额则分别在原子、留数与 M\"obius--Abel 三层被完全显式记账。

\endinput
